\documentclass[10pt,a4paper]{amsart}
\usepackage[margin=19mm]{geometry}
\usepackage{amsmath,amssymb,mathtools}
\usepackage[scr=rsfs,cal=euler]{mathalfa}
\usepackage{xfrac}
\usepackage[unicode,hidelinks,bookmarks=false]{hyperref}
\newcommand{\ie}{i.e.\ }
\newcommand{\cf}{cf.\ }
\newcommand{\resp}{respectively\ }
\newcommand{\Subsection}{\subsection}
\providecommand{\nobreakdash}{\nobreak\hbox{-}}
\setlength{\emergencystretch}{2em}
\multlinegap0pt
\usepackage{bm}
\usepackage{bbm}
\usepackage{dsfont}
\usepackage{xspace}
\usepackage{cancel}
\usepackage{mathtools}
\usepackage[shortlabels]{enumitem}
\usepackage{amsthm}
\makeatletter
\@ifundefined{theorem}{\newtheorem{theorem}{Theorem}}{}
\@ifundefined{lemma}{\newtheorem{lemma}[theorem]{Lemma}}{}
\makeatother
\newcommand{\norm}[1]{\lVert#1\rVert}
\title[Existence and multiplicity of normalized solutions to a larg]{Existence and multiplicity of normalized solutions to a large class of elliptic equations on bounded domains with general boundary conditions}
\author{Claudianor O. Alves, Zhentao He, Chao Ji}
\date{Source: arXiv:2507.04624v1 (CC BY 4.0)}
\begin{document}
\maketitle

\noindent\emph{Source and licence.} Existence and multiplicity of normalized solutions to a large class of elliptic equations on bounded domains with general boundary conditions. Version arXiv:2507.04624v1 (CC BY 4.0), distributed under the Creative Commons Attribution 4.0 International licence, as recorded in the arXiv licence metadata for this exact version and shown on the abstract page https://arxiv.org/abs/2507.04624v1. Full rights notice: licenses/arxiv-2507.04624-CC-BY-4.0.txt.\par\smallskip

\noindent\textit{Editorial scope.} Counted unit: the Lemma whose source label is \texttt{lemnex} in the article (source0.tex lines 782--794, source label \texttt{lemnex}), delivered together with the article's own complete original proof of it (source0.tex lines 796--847, one \texttt{proof} environment of 52 lines). No mathematics is written, repaired or re-derived: statement and proof are the article's own text line for line. Required context retained uncounted: eqgn (lines 503--505); lembounded (lines 594--596). Every definition, display and label of those lines is retained unchanged. Omitted from this package and not redistributed: 1--502, 506--593, 597--781, 795--795, 848--1944. Nothing inside a retained excerpt is omitted or altered: every retained body is delivered exactly as the article's lines are written, so no omission receipt of any kind accompanies this package. Citations: the retained excerpts carry no citation command at all, so no bibliography entry of the article is transcribed and no reference is redistributed. Reference adaptation: none was needed and none was made; every reference inside the delivered unit resolves inside it, so no unresolved reference or citation is produced. Printed slips kept literal: no printed word, symbol, hypothesis or number of the counted statement or of its proof was altered. Layout and class: the article's own class is \texttt{amsart} with packages including amsmath, amssymb, amsthm, latexsym, cite, cancel, caption, graphicx, wasysym, overpic, tikz, color, subfigure, hyperref; the wrapper is the lane's pinned \texttt{amsart} editorial wrapper at 10pt on A4, which restates the theorem-like environments the retained text needs and adds the article's own macro definitions that the retained text needs (\texttt{\textbackslash norm}), with extra wrapper packages (bm, bbm, dsfont, xspace, cancel, mathtools, enumitem). The delivered numbering is the unit's own. Assets: no retained span contains a figure environment, a graphics-inclusion command, a PDF-page inclusion, a table or a TikZ picture; no image file is copied into this package, the package declares no asset dependency and no image file is redistributed with it. Licence: version arXiv:2507.04624v1 of this article is distributed under the Creative Commons Attribution 4.0 International licence, as recorded in the arXiv licence metadata for this exact version and shown on the abstract page https://arxiv.org/abs/2507.04624v1. A full rights notice is delivered at licenses/arxiv-2507.04624-CC-BY-4.0.txt. Characters: every retained excerpt is pure ASCII, so the delivered engine, XeLaTeX with its default Latin Modern text font, reports no missing character.
\begin{equation}\label{eqgn}
    \|u\|_{p}^p \leq C_{p,N} \|u\|_{2}^{(1 - \beta_p)p} \|\nabla u\|_{2}^{\beta_p p} \quad \forall u \in  H_0^1(\Omega),
\end{equation}

\begin{lemma}\label{lembounded}
    For $r>1$ sufficiently large, $\{u_{n,r}\}_{n\geq 1}$ is bounded in $H^1_0(\Omega)$.
\end{lemma}

\begin{lemma}\label{lemnex}
The following two non-existence results hold:
    \begin{enumerate}[label=\rm(\roman*)]
        \item Assume that ($f_1$)-($f_3$) hold and that $\mu \in (0,1)$. Then, given $M>0$ independent of $\mu \in (0,1)$ and $r>0$, there exists a constant $\mu_0:=\mu_0(\Omega,f,M) \in (0,1)$ such that, for any $0<\nu<\mu_0$, there exists no $u \in H^1_0(\Omega)$ satisfying
             $$
             E'(u) = 0,\,\,E(u)\leq M,\,\,\text{and}\,\,\int_{\Omega}|u|^2\,dx=\nu.
             $$
        \item Assume that ($f_1$) and ($f_4$) hold. Then, for any $M>0$  independent of $\mu >0$ and $0<\nu<\mu^*$, there exists no $u \in H_0^1(\Omega)$ satisfying
             $$
             E'(u) = 0,\,\,E(u)\leq M\mu^*\,\,\text{and}\,\,\int_{\Omega}|u|^2 \,dx=\nu.
             $$
    \end{enumerate}
\end{lemma}

\begin{proof}
    (i) Suppose that ($f_1$)-($f_3$) hold and that $\mu \in (0,1)$. For any fixed $M>0$, we claim that, there exists $R>0$ independent of $\mu \in (0,1)$ and $r>1$, such that for all  $u \in H_0^1(\Omega)$ satisfying $E'(u) = 0$ and $E(u)\leq M$, we have $\norm{u}\leq R$. Indeed, if not, then there exists a sequence $\{u_n\}\subset H_0^1(\Omega)$ satisfying $E'(u_n) = 0$, $E(u_n)\leq M$
    and $\norm{u_n} \to +\infty$ as $n \to +\infty$. Repeating the arguments of Lemma \ref{lembounded}, we obtain a contradiction.

    Now, argue by contradiction, assume that for some $\nu>0$, there exists $u \in H_0^1(\Omega)$ satisfying $E'(u) = 0$, $E(u)\leq M$ and $\int_{\Omega}|u|^2\,dx=\nu$. By $E'(u) = 0$, ($f_1$) and \eqref{eqgn}, one has
    \begin{equation*}\label{lemnonex}
        \begin{aligned}
        \norm{u}^2 = \int_\Omega f(u)u\,dx &\leq K_2 \norm{u}_2^2+K_{p} C_{p,N}\norm{u}^{\beta_{p}{p}}\norm{u}_2^{(1-\beta_{p}){p}} \\
        &\leq \frac{K_2}{\lambda_{1,0,0}} \norm{u}^2+K_{p} C_{p,N}\norm{u}^{\beta_{p}{p}}\norm{u}_2^{(1-\beta_{p}){p}}.
        \end{aligned}
    \end{equation*}
    Hence,
    \begin{equation*}
    \frac{\lambda_{1,0,0}-K_2}{\lambda_{1,0,0}}\leq K_{p} C_{p,N}\norm{u}^{\beta_{p}{p}-2}\nu^\frac{(1-\beta_{p}){p}}{2}.
    \end{equation*}
    If $p \leq 2+\frac{4}{N}$, then $\beta_pp\leq 2$ and
    $$
    \frac{\lambda_{1,0,0}-K_2}{\lambda_{1,0,0}}\leq K_{p} C_{p,N}\lambda_{1,0,0}^\frac{{\beta_{p}{p}-2}}{2}\nu^\frac{{\beta_{p}{p}-2}}{2} \nu^\frac{(1-\beta_{p}){p}}{2}=K_{p} C_{p,N} \lambda_{1,0,0}^\frac{{\beta_{p}{p}-2}}{2}\nu^\frac{p-2}{2}.
    $$
     If $p > 2+\frac{4}{N}$, then $\beta_pp > 2$ and
    $$
    \frac{\lambda_{1,0,0}-K_2}{\lambda_{1,0,0}}\leq K_{p} C_{p,N}R^{\beta_{p}{p}-2} \nu^\frac{(1-\beta_{p}){p}}{2}.
    $$
    Therefore, there exists $\mu_0>0$ such that, for any $0<\nu<\mu_0$, there exists no $u \in H_0^1(\Omega)$ satisfying
     $$
     E'(u) = 0,\,\,E(u)\leq M\,\,\text{and}\,\,\int_{\Omega}|u|^2 \,dx=\nu.
     $$


    (ii) Suppose that ($f_1$) and ($f_4$) hold. Given $M>0$, argue by contradiction, assume that there exists $u \in H_0^1(\Omega)$ satisfying $E'(u) = 0$, $E(u)\leq M\mu^*$ and $\int_{\Omega}|u|^2\,dx=\nu$ for some $0<\nu<\mu^*$. Then, by ($f_4$), we have
    $$
    \norm{u}^2 \leq \frac{2qE(u)-2\langle E'(u),u\rangle}{q-2}\leq \frac{2qM\mu^*}{q-2}.
    $$
    Repeating the previous argument, we derive that, if $2<p\leq2+\frac{4}{N}$, then
    $$
    \frac{\lambda_{1,0,0}-K_2}{\lambda_{1,0,0}}\leq K_{p} C_{p,N} \lambda_{1,0,0}^\frac{{\beta_{p}{p}-2}}{2}\nu^\frac{p-2}{2},
    $$
    and, if $2+\frac{4}{N}<p<2^*$, then
    $$
    \begin{aligned}
        \frac{\lambda_{1,0,0}-K_2}{\lambda_{1,0,0}}&\leq K_{p} C_{p,N}\norm{u}^{\beta_{p}{p}-2}\nu^\frac{(1-\beta_{p}){p}}{2}\\
        &\leq K_{p} C_{p,N}\left(\frac{2qM\mu^*}{q-2}\right)^\frac{{\beta_{p}{p}-2}}{2}{\nu}^\frac{(1-\beta_p) p}{2}\\
        &<K_{p} C_{p,N}\left(\frac{2qM}{q-2}\right)^\frac{{\beta_{p}{p}-2}}{2}{\mu^*}^\frac{p-2}{2},
    \end{aligned}
    $$
    which leads to a contradiction, since $\frac{2-\beta_pp}{p-2}=\frac{2}{p-2}-\frac{N}{2}$ and
 $$\nu < \mu^*=
\begin{cases}
    \left(\frac{\lambda_{1,0,0}-K_2}{K_pC_{p,N}}\right)^{\frac{2}{p-2}}\lambda_{1,0,0}^{-\frac{N}{2}}  &\text{ if } 2<p\leq 2+\frac{4}{N},\\
    \left(\frac{\lambda_{1,0,0}-K_2}{K_pC_{p,N}\lambda_{1,0,0}}\right)^{\frac{2}{p-2}}\left(\frac{2qM}{q-2}\right)^{\frac{2}{p-2}-\frac{N}{2}} &\text{ if }  2+\frac{4}{N}<p<2^*.
\end{cases}$$
\end{proof}

\end{document}
